\section{Modular cancellation and the scale of a collision}\label{sec:collisions}

The positive quantity $C=\sum_{n\ge2}H_n$ is not the complete field.
In $H=H_1+C$, modular reflection forces cancellations that are invisible
in a separate bound for either term. We first make those cancellations
explicit and then examine their effect at a hypothetical double zero.

\subsection{The endpoint expansion of a cosine transform}

\begin{lemma}\label{lem:endpoint}
Suppose $g$ is smooth on $[0,\infty)$ and
$\sup_{u\ge0}(1+u)^k|g^{(j)}(u)|<\infty$ for all integers $j,k\ge0$.
Then every derivative is integrable, all the boundary terms below vanish
at infinity, and,
for each $M\ge1$ and $x\ne0$,
\begin{equation}\label{eq:endpoint-exact}
 \int_0^\infty g(u)\cos(xu)\dd u
 =\sum_{k=1}^M\frac{(-1)^k g^{(2k-1)}(0)}{x^{2k}}
   +\frac{(-1)^M}{x^{2M}}
       \int_0^\infty g^{(2M)}(u)\cos(xu)\dd u.
\end{equation}
If the first nonzero odd endpoint derivative has order $p-1$, where
$p=2r$, and $A=(-1)^r g^{(p-1)}(0)\ne0$, then its transform $V$ satisfies
\begin{align}
 V(x)&=Ax^{-p}+O(x^{-p-2}),\notag\\
 V'(x)&=-pAx^{-p-1}+O(x^{-p-3}),\label{eq:endpoint-jets}\\
 V''(x)&=p(p+1)Ax^{-p-2}+O(x^{-p-4}).\notag
\end{align}
\end{lemma}
\begin{proof}
One pair of integrations gives
$\int g\cos=-g'(0)/x^2-x^{-2}\int g''\cos$.
Iterating this equality proves \eqref{eq:endpoint-exact}. Taking one extra
pair bounds the remainder after the $M$th coefficient by a constant times
$x^{-2M-2}$.

For the derivative assertions, start with the differentiated integrals
rather than differentiating a big-$O$ term. An integration by parts gives
\[
 V'(x)=-\int_0^\infty ug(u)\sin(xu)\dd u
       =-\frac1x\int_0^\infty (ug(u))'\cos(xu)\dd u.
\]
The $k$th endpoint derivative of $(ug)'$ is $(k+1)g^{(k)}(0)$.
Thus the leading coefficient in this cosine transform is $pA$.
Similarly, $V''=-\int u^2g\cos$, and
$(u^2g)^{(k)}(0)=k(k-1)g^{(k-2)}(0)$ for $k\ge2$.
Its first nonzero odd endpoint has order $p+1$ and coefficient
$-p(p+1)A$ in the cosine expansion. Applying \eqref{eq:endpoint-exact}
to these two functions proves all three estimates.
\end{proof}

\subsection{The full endpoint hierarchy}

The evenness of $w_t=e^{tu^2}\Ph$ proved in Section~\ref{sec:theta} gives
$w_t^{(2k-1)}(0)=0$ for every $k\ge1$. Write
\[
 w_t=g_{1,t}+g_{A,t}+g_{R,t},
\]
where the three terms are the weighted sources of $H_1$, $A_{\mathcal P}$
and $R_{\mathcal P}$. The complete modular identity therefore gives
\begin{equation}\label{eq:all-endpoints}
 g_{1,t}^{(2k-1)}(0)+g_{A,t}^{(2k-1)}(0)
                         +g_{R,t}^{(2k-1)}(0)=0
 \qquad(k\ge1).
\end{equation}
For example, multiplication by the heat weight changes the third source
derivative to $\phi'''(0)+6t\phi'(0)$, but the sum still vanishes.
No finite list of endpoint checks is being substituted for
\eqref{eq:inversion}.

Let
\begin{equation}\label{eq:exact-masses}
 b=\phi_1'(0),\qquad q_A=-g_{A,t}'(0),\qquad
 q_R=-g_{R,t}'(0).
\end{equation}
For a nonempty finite prime set, $q_A,q_R>0$ by
Theorem~\ref{thm:positive}, and these quantities do not depend on $t$.
The first equation in \eqref{eq:all-endpoints} says
\begin{equation}\label{eq:mass-cancellation}
 b=q_A+q_R=\sum_{n\ge2}-\phi_n'(0)>0.
\end{equation}
In particular, the first mode has the opposite endpoint sign from the
higher modes. Lemma~\ref{lem:endpoint} gives, at every fixed $t$ in the
time interval,
\begin{equation}\label{eq:three-tails}
 H_1=-\frac b{x^2}+O(x^{-4}),\qquad
 A_{\mathcal P}=\frac{q_A}{x^2}+O(x^{-4}),\qquad
 R_{\mathcal P}=\frac{q_R}{x^2}+O(x^{-4}).
\end{equation}
Every algebraic coefficient cancels in their complete sum by
\eqref{eq:all-endpoints}. At the first level this is the exact identity
\begin{equation}\label{eq:coupled-second}
 H(t,x)=-\frac1{x^2}\int_0^\infty
   \bigl(g_{1,t}''+g_{A,t}''+g_{R,t}''\bigr)\cos(xu)\dd u.
\end{equation}
The first summand in this differentiated source has not acquired a
nonnegative sign. Thus \eqref{eq:coupled-second} preserves the cancellation
but is not a positivity formula for $H$.

\begin{proposition}[A fixed omitted tail changes the asymptotics]\label{prop:omission}
For each fixed nonempty finite prime set and each fixed
$0\le t\le1/2$,
\[
 H_1(t,x)+A_{\mathcal P}(t,x)
   =-\frac{q_R}{x^2}+O(x^{-4})<0
\]
for all sufficiently large $|x|$. Moreover,
\begin{equation}\label{eq:ratio-limits}
 \frac{R_{\mathcal P}}{A_{\mathcal P}}\longrightarrow\frac{q_R}{q_A}>0,
 \qquad
 1+\frac{H_1}{A_{\mathcal P}}\longrightarrow-\frac{q_R}{q_A}.
\end{equation}
\end{proposition}
\begin{proof}
Use $H_1+A_{\mathcal P}=H-R_{\mathcal P}$, the rapid decay of $H$,
and \eqref{eq:three-tails}. Division by the positive leading term of
$A_{\mathcal P}$ gives the two limits.
\end{proof}

An accurate relative approximation to the positive arithmetic tail is
therefore not a relative approximation to the complete field near its
zeros. Proposition~\ref{prop:omission} concerns omission of a fixed
finite-prime remainder; it says nothing against an infinite assembly with
its own convergence and error proof.

\subsection{A regular normalization and the signed identity}

Assume now that $2\in\mathcal P$, and suppress this subscript temporarily.
The positive denominator in \eqref{eq:lower} defines, on the real axis,
\begin{equation}\label{eq:normalized}
 h=\frac HA,\qquad Z=1+\frac{H_1}A,\qquad r=\frac RA,
 \qquad h=Z+r,\qquad \alpha=\frac{A_x}{A}.
\end{equation}
These functions are holomorphic wherever the carrier is nonzero.
The strip positivity in Section~\ref{sec:positivity} supplies such a
domain containing the classical strip $|\Im z|\le1$. On the real axis
all these functions are real analytic.
The carrier $A$ is specified by the source, independently of the value
of $H$ or the position of a possible collision.
Because both $H$ and $A$ satisfy the heat equation,
\[
 A_t h+A h_t=-(A_{xx}h+2A_xh_x+A h_{xx}),\qquad A_t=-A_{xx},
\]
so
\begin{equation}\label{eq:drift}
 h_t=-h_{xx}-2\alpha h_x.
\end{equation}

Define the physical energy and the first signed Laguerre expression by
\begin{equation}\label{eq:energy}
 E=H^2+H_x^2,\qquad \Lap_0=H_x^2-HH_{xx}.
\end{equation}
From $H_x=A(h_x+\alpha h)$ and
$H_{xx}=A(h_{xx}+2\alpha h_x+(\alpha_x+\alpha^2)h)$, we obtain
\begin{equation}\label{eq:laguerre-normalized}
 \boxed{\Lap_0=A^2(h_x^2-hh_{xx})-\alpha_xH^2.}
\end{equation}
The cross terms cancel before estimation. Theorem~\ref{thm:jets} gives
the definite bound
\begin{equation}\label{eq:payment}
 |\alpha_x H^2|\le H^2/64\le E/64.
\end{equation}
This is an energy-scale estimate for the carrier term. The remaining
quantity $A^2(h_x^2-hh_{xx})$ has no sign established here.
The identity alone does not reduce the strength of the missing estimate.

\subsection{What changing the carrier can accomplish}

Away from zeros of $H$, logarithmic differentiation of $h=H/A$ gives
\begin{equation}\label{eq:carrier-invariance}
 \left(\frac{h_x}{h}\right)^2-\frac{h_{xx}}h
 =\frac{\Lap_0}{H^2}+(\log A)_{xx}.
\end{equation}
For two positive carriers $A,B$, the corresponding expressions differ
by $(\log(A/B))_{xx}$. At a stationary point of $H$, the first term on the
right is $-H_{xx}/H$. Thus a bounded change of carrier curvature changes
only a bounded part of this quantity; it cannot remove the singular
collision behavior calculated below.

One may also use the full positive carrier $C=\sum_{n\ge2}H_n$.
The strip and logarithmic-derivative estimates apply to $C$ with $m=2$,
and $H/C=1+H_1/C$ contains no omitted arithmetic remainder. The cancellation
with the first mode is still present. This limits what carrier optimization
alone establishes; it does not exclude estimates using the complete modular
relation quantitatively.

\subsection{Local behavior at a hypothetical double event}

We give the local calculation explicitly so that the scale in
\eqref{eq:payment} can be assessed. Suppose a real point $(T,x_*)$ satisfies
\begin{equation}\label{eq:double}
 H(T,x_*)=H_x(T,x_*)=0,\qquad c=\tfrac12H_{xx}(T,x_*)\ne0.
\end{equation}
For the packet bounds we restrict $0<T\le1/2$, although the following
local heat calculation is more general. Existence of such an event for
the theta field is not asserted.

Let $s=T-t>0$ and $x=x_*+\sqrt{s}\,y$. For a smooth heat solution $g$,
write $g_*^{(j)}=\partial_x^jg(T,x_*)$, and suppose $|y|\le R$.
Let $M_j(g)$ bound $|\partial_x^jg|$ on a rectangle containing the
time and spatial segments in the Taylor expansions below.
Taylor expansion in time gives
\[
 g(T-s,x)=g(T,x)+s g_{xx}(T,x)+\mathcal E_t,\qquad
 |\mathcal E_t|\le\tfrac12s^2M_4(g).
\]
Expand the first term through its second spatial derivative and the
second through its constant spatial term. Their two remainders are
bounded respectively by $s^{3/2}R^3M_3/6$ and $s^{3/2}RM_3$.
Hence
\begin{align}
 g(T-s,x_*+\sqrt{s}\,y)
 &=g_*+\sqrt{s}\,y\,g_*'
           +s(1+y^2/2)g_*''+\mathcal R_g,\label{eq:local}\\
 |\mathcal R_g|
 &\le s^{3/2}(R+R^3/6)M_3(g)+\tfrac12s^2M_4(g).\label{eq:local-error}
\end{align}
Applying the same argument to $g_x$ shifts each derivative index by one.
For $H$, \eqref{eq:double} gives
\begin{equation}\label{eq:collision-limit}
 \left|\frac{H(T-s,x_*+\sqrt{s}\,y)}{cs}-(y^2+2)\right|
 \le\frac{\sqrt{s}(R+R^3/6)M_3(H)+sM_4(H)/2}{|c|}.
\end{equation}
In particular $H/(cs)\to y^2+2$ and $H_x/(c\sqrt{s})\to2y$
uniformly on bounded $y$ intervals. Compactness of $(T,x_*)$ alone
does not remove the explicit factor $1/|c|$.

The implicit function theorem, applied to $H_x(T-s,x)=0$, gives a
real analytic stationary branch $x_{\rm stat}(s)=x_*+O(s)$.
Its derivative at zero is $-H_{xxx}(T,x_*)/(2c)$.
Substitution into the Taylor formula yields
\begin{align}
 H(T-s,x_{\rm stat}(s))&=2cs+O(s^2),&
 H_{xx}(T-s,x_{\rm stat}(s))&=2c+O(s),\notag\\
 E&=4c^2s^2+O(s^3),&
 \Lap_0&=-4c^2s+O(s^2).\label{eq:stationary}
\end{align}
Thus
\begin{equation}\label{eq:singular-ratio}
 \frac{\Lap_0}{E}=-\frac1s+O(1).
\end{equation}
This calculation explains why a bound by a fixed multiple of $E$ is
materially different from a small absolute error or a bound proportional
to $|H|$.

\subsection{Matching the complete exterior packet}

At an event \eqref{eq:double}, decomposition \eqref{eq:decomposition}
and its first two derivatives give
\begin{equation}\label{eq:event-matching}
 A_*=-H_{1,*}-R_*,\quad
 A_{x,*}=-H_{1,x,*}-R_{x,*},\quad
 A_{xx,*}=2c-H_{1,xx,*}-R_{xx,*}.
\end{equation}
For example, the exterior combination $W=H_1A_x-H_{1,x}A$ obeys
\[
 W_*=H_{1,x,*}R_*-H_{1,*}R_{x,*},\qquad
 |W_*|\le |H_{1,x,*}|\epsilon_0(m)+|H_{1,*}|\epsilon_1(m).
\]
At an arbitrary stationary point $H_x=0$, the exact formula is
$W=-H_{1,x}H+H_{1,x}R-H_1R_x$. These are necessary matching equations,
not independently signed restrictions.

The actual exterior arguments inside $A$ remain those of
\eqref{eq:complete-packet}. At $t=T-s$ their displacement is
$i(T-s)\log d$, so the rescaled coordinate is
$y+i(T-s)\log d/\sqrt{s}$. For $d>1$, this escapes every bounded local
coordinate domain. A Taylor polynomial valid for bounded $y$ cannot
evaluate these exterior observations without a separate error argument.
Equations \eqref{eq:local}--\eqref{eq:local-error} instead expand the
\emph{whole packet in its base point}, retaining its exterior data.

In the normalized notation \eqref{eq:normalized}, event matching is
\begin{equation}\label{eq:normalized-matching}
 Z_*=-r_*\in[-\rho_m,0),\qquad
 Z_{x,*}=-r_{x,*},\qquad
 Z_{xx,*}+r_{xx,*}=\frac{2c}{A_*}.
\end{equation}
Theorem~\ref{thm:jets} supplies uniform tolerances for the first two
equalities. It gives no sign for the third. Moreover,
\eqref{eq:ratio-limits} shows that the narrow value band is approached at
large height even in the absence of a collision.

Along the stationary predecessor of the event,
$h=2cs/A_*+O(s^2)$, whereas $r$ tends to the strictly positive number
$R_*/A_*$. Thus a fixed error bound for $r$, however small, is not an
error of order $h^2$. The bounded term \eqref{eq:payment} has been
controlled; the complete signed residual in \eqref{eq:laguerre-normalized}
has not. Establishing a collision exclusion would require an additional
estimate on that residual with the modular cancellation retained.
